\section{Conclusion}
System1Bench measures reference agreement, joint correctness, correctness turnover, prediction stability and operational efficiency at a fixed typed decision interface. The source atlas and independent legal/scientific diagnostics keep workload type and reference provenance visible instead of turning unlike labels into a pooled score. On SciFact's balanced three-way cited-pair diagnostic, Qwen's near-unchanged total accuracy under candidate reversal conceals a large loss on no-information cases offset by gains on support and contradiction cases. A fresh post-hoc census of all eligible cited pairs shows the same class-level attrition outside that selected sample, while the aggregate difference remains imprecise; the derived no-information reference is not an adjudicated neutral class. The legal four-cell control goes further: superficially similar reversal outcomes arise from different couplings to displayed candidate order and answer-code identity in the two 8B adapters, with exact-score ties requiring their own sensitivity check. The supervised, document-blind ContractNLI prior and the mixed held-out wording test illustrate why source shortcuts and attractive post-hoc mechanisms must also be tested. Together these results suggest a practical reporting rule: preserve source-specific references, paired class transitions, orthogonal interface controls and the execution boundary. They characterize downstream decision behavior, not whole-system completion or deployment utility; those claims require evaluation of the full workflow.
